\hypertarget{python-2.2-to-2.3-type-class-unification-descriptors-c3-mro}{%
\chapter{Python 2.2 to 2.3: Type-Class Unification, Descriptors, \& C3
MRO}\label{python-2.2-to-2.3-type-class-unification-descriptors-c3-mro}}

\hypertarget{type-class-unification-the-inception-of-new-style-classes-pep-252-pep-253}{%
\subsection{4.1 Type-Class Unification: The Inception of New-Style
Classes (PEP 252 \& PEP
253)}\label{type-class-unification-the-inception-of-new-style-classes-pep-252-pep-253}}

Before Python 2.2, user-defined classes and built-in types existed as
entirely separate entities under the hood. This architectural separation
introduced severe language inconsistencies.

\hypertarget{the-classic-class-inconsistency}{%
\subsubsection{1. The Classic Class
Inconsistency}\label{the-classic-class-inconsistency}}

In the classic class model (pre-Python 2.2): * \textbf{Classic Classes}:
Created using the \passthrough{\lstinline!class MyClass:!} syntax
without inheriting from any class. All user-defined classes were
represented by a single C type, \passthrough{\lstinline!PyClass\_Type!},
defined by the \passthrough{\lstinline!PyClassObject!} struct. *
\textbf{Classic Instances}: Every instance of any user-defined class was
represented by a single C type,
\passthrough{\lstinline!PyInstance\_Type!}, mapped to a
\passthrough{\lstinline!PyInstanceObject!} struct. * \textbf{The
Dichotomy}: The \passthrough{\lstinline!type()!} of any classic class
instance always returned \passthrough{\lstinline!<type 'instance'>!},
rather than the class itself. The actual class of the instance was
accessible only via the \passthrough{\lstinline!\_\_class\_\_!}
attribute: ```python \# Python 1.5 - 2.1 Classic Class Behavior class
Classic: pass

c = Classic() print(type(c)) \# \textless type `instance'\textgreater{}
print(c.\_\_class\_\_) \# \textless class
`\textbf{main}.Classic'\textgreater{}
``\passthrough{\lstinline!* **Built-in Types**: In contrast, built-in types (e.g.,!}int\passthrough{\lstinline!,!}list\passthrough{\lstinline!,!}dict\passthrough{\lstinline!,!}str\passthrough{\lstinline!) were instances of!}PyType\_Type\passthrough{\lstinline!, represented in C by a static!}PyTypeObject\passthrough{\lstinline!struct. * **The Subclassing Barrier**: Because of this dichotomy, user-defined classes could not inherit from built-in types. Subclassing!}list\passthrough{\lstinline!or!}dict`
to extend their behavior was impossible, because the classic class
object engine could not parse or manage the memory layouts of built-in C
types.

\hypertarget{the-solution-new-style-classes-and-object}{%
\subsubsection{\texorpdfstring{2. The Solution: New-Style Classes and
\texttt{object}}{2. The Solution: New-Style Classes and object}}\label{the-solution-new-style-classes-and-object}}

Introduced in Python 2.2 by PEP 252 and PEP 253, \textbf{new-style
classes} unified types and classes. * \textbf{The Unified Base}: A
new-style class is defined by inheriting (directly or indirectly) from
the built-in \passthrough{\lstinline!object!} type (the root of the
unified type tree). * \textbf{Type Unification}: Under the unified
model, user-defined classes are themselves type objects (instances of
\passthrough{\lstinline!type!}), just like built-in types. Checking
\passthrough{\lstinline!type(c)!} now returns the class itself.
```python \# Python 2.2+ New-Style Class Behavior class
NewStyle(object): pass

n = NewStyle() print(type(n)) \# \textless class
`\textbf{main}.NewStyle'\textgreater{} ```

\hypertarget{low-level-c-representation-pytypeobject}{%
\subsubsection{\texorpdfstring{3. Low-Level C Representation:
\texttt{PyTypeObject}}{3. Low-Level C Representation: PyTypeObject}}\label{low-level-c-representation-pytypeobject}}

Every class (built-in or new-style user-defined) is represented at the
C-level as an instance of \passthrough{\lstinline!PyTypeObject!}. Let's
inspect the core fields in CPython's
\passthrough{\lstinline!Include/object.h!}:

\begin{lstlisting}[language=C]
typedef struct _typeobject {
    PyObject_VAR_HEAD
    const char *tp_name;                 /* For printing, in format "<module>.<name>" */
    Py_ssize_t tp_basicsize, tp_itemsize; /* For allocation sizes */

    /* Methods to implement standard operations */
    destructor tp_dealloc;
    printfunc tp_print;
    getattrfunc tp_getattr;
    setattrfunc tp_setattr;
    
    /* Attribute lookup slot */
    getattrofunc tp_getattro;            /* Pointing to PyObject_GenericGetAttr */
    setattrofunc tp_setattro;            /* Pointing to PyObject_GenericSetAttr */

    /* Protocol slot mappings */
    PyNumberMethods *tp_as_number;
    PySequenceMethods *tp_as_sequence;
    PyMappingMethods *tp_as_mapping;

    /* Inheritance and lookup structures */
    PyObject *tp_dict;                  /* Namespace dictionary */
    PyObject *tp_bases;                 /* Tuple of base classes */
    PyObject *tp_mro;                   /* Method Resolution Order tuple */
    
    /* Allocation / Initialization slots */
    newfunc tp_new;                     /* __new__ allocation entry */
    initproc tp_init;                   /* __init__ initialization entry */
    allocfunc tp_alloc;                 /* Low-level memory allocator */
    
    /* Flags and inheritance info */
    unsigned long tp_flags;
    struct _typeobject *tp_base;        /* Direct base pointer */
} PyTypeObject;
\end{lstlisting}

\hypertarget{slot-wrapper-descriptors}{%
\subsubsection{4. Slot Wrapper
Descriptors}\label{slot-wrapper-descriptors}}

Because C slots (like \passthrough{\lstinline!tp\_init!}) expect
standard C function signatures (e.g.,
\passthrough{\lstinline!int (*tp\_init)(PyObject *, PyObject *, PyObject *)!}),
but user-defined classes write Python methods
(\passthrough{\lstinline!def \_\_init\_\_(self, ...)!}), CPython
utilizes \textbf{slot wrappers} to bridge the boundary: * \textbf{Python
to C (Slot Fillers)}: When a new-style class defines a Python method
like \passthrough{\lstinline!\_\_init\_\_!}, the type compiler
dynamically wraps this method in a C function (such as
\passthrough{\lstinline!slot\_tp\_init!}) and assigns it to the type
object's \passthrough{\lstinline!tp\_init!} slot. When
\passthrough{\lstinline!tp\_init!} is called,
\passthrough{\lstinline!slot\_tp\_init!} executes the Python function. *
\textbf{C to Python (Wrapper Descriptors)}: To expose built-in C slots
(like \passthrough{\lstinline!object!}'s default initialization code) to
Python, CPython wraps them in descriptor objects
(e.g.~\passthrough{\lstinline!\_\_init\_\_!} is exposed as
\passthrough{\lstinline!<slot wrapper '\_\_init\_\_' of 'object' objects>!}).

\begin{center}\rule{0.5\linewidth}{0.5pt}\end{center}

\hypertarget{the-descriptor-protocol-attribute-lookup-chain}{%
\subsection{4.2 The Descriptor Protocol \& Attribute Lookup
Chain}\label{the-descriptor-protocol-attribute-lookup-chain}}

The descriptor protocol is the underlying mechanism that enables
properties, methods, classmethods, and staticmethods in Python.

\hypertarget{the-descriptor-protocol-definition}{%
\subsubsection{1. The Descriptor Protocol
Definition}\label{the-descriptor-protocol-definition}}

A descriptor is an object that implements at least one of the three
descriptor protocol methods:

\begin{lstlisting}[language=Python]
def __get__(self, instance, owner=None):
    """Invoked when getting the attribute."""
    pass

def __set__(self, instance, value):
    """Invoked when setting the attribute."""
    pass

def __delete__(self, instance):
    """Invoked when deleting the attribute."""
    pass
\end{lstlisting}

\begin{itemize}
\tightlist
\item
  \textbf{Data Descriptor}: Implements both
  \passthrough{\lstinline!\_\_get\_\_!} AND
  \passthrough{\lstinline!\_\_set\_\_!} (and/or
  \passthrough{\lstinline!\_\_delete\_\_!}).
\item
  \textbf{Non-Data Descriptor}: Implements only
  \passthrough{\lstinline!\_\_get\_\_!} (e.g.~functions, classmethods,
  staticmethods).
\end{itemize}

\hypertarget{the-attribute-lookup-algorithm-pyobject_genericgetattr}{%
\subsubsection{\texorpdfstring{2. The Attribute Lookup Algorithm
(\texttt{PyObject\_GenericGetAttr})}{2. The Attribute Lookup Algorithm (PyObject\_GenericGetAttr)}}\label{the-attribute-lookup-algorithm-pyobject_genericgetattr}}

When an attribute is accessed (\passthrough{\lstinline!obj.name!}) on a
new-style instance \passthrough{\lstinline!obj!} of class
\passthrough{\lstinline!Class!}, CPython invokes the
\passthrough{\lstinline!tp\_getattro!} slot of the type object. By
default, this points to
\passthrough{\lstinline!PyObject\_GenericGetAttr!} in
\passthrough{\lstinline!Objects/object.c!}.

The lookup follows this strict resolution flow: 1. \textbf{MRO Search}:
CPython calls \passthrough{\lstinline!\_PyType\_Lookup(Class, name)!}.
This searches the class's namespace dictionary
\passthrough{\lstinline!tp\_dict!} and the namespace dictionaries of all
parent classes in the Method Resolution Order
(\passthrough{\lstinline!tp\_mro!}). Let the resolved object be
\passthrough{\lstinline!descr!}. 2. \textbf{Data Descriptor Check}: If
\passthrough{\lstinline!descr!} is found, and its type implements the
\passthrough{\lstinline!\_\_set\_\_!} (or
\passthrough{\lstinline!\_\_delete\_\_!}) slot
(\passthrough{\lstinline!descr->ob\_type->tp\_descr\_set!} is not NULL):
* Call
\passthrough{\lstinline!descr->ob\_type->tp\_descr\_get(descr, obj, Class)!}
and return the result immediately. 3. \textbf{Instance Dict Lookup}:
Check the instance dictionary of \passthrough{\lstinline!obj!} (the
\passthrough{\lstinline!\_\_dict\_\_!} table at
\passthrough{\lstinline!obj->ob\_dict!}). If
\passthrough{\lstinline!name!} exists in the instance dictionary, return
its associated value. 4. \textbf{Non-Data Descriptor Check}: If
\passthrough{\lstinline!descr!} is found: * If
\passthrough{\lstinline!descr!} has a
\passthrough{\lstinline!\_\_get\_\_!} slot
(\passthrough{\lstinline!descr->ob\_type->tp\_descr\_get!} is not NULL):
* Call
\passthrough{\lstinline!descr->ob\_type->tp\_descr\_get(descr, obj, Class)!}
and return the result. * If \passthrough{\lstinline!descr!} does not
have a \passthrough{\lstinline!\_\_get\_\_!} slot (e.g.~a plain class
attribute like a string or int): * Return
\passthrough{\lstinline!descr!} directly. 5. \textbf{Fallback to
\passthrough{\lstinline!\_\_getattr\_\_!}}: If the attribute is still
not resolved, CPython raises an
\passthrough{\lstinline!AttributeError!}, which triggers the invocation
of the fallback \passthrough{\lstinline!\_\_getattr\_\_(self, name)!}
method if defined on the class.

\begin{lstlisting}
       [Start Lookup: obj.name]
                  |
                  v
       [Search MRO for "name"]
                  |
       +----------+----------+
       |                     |
   (Not Found)            (Found)
       |                     |
       |                     v
       |          [Is it a Data Descriptor?]
       |          (has __get__ and __set__)
       |                     |
       |             +-------+-------+
       |             | Yes           | No
       |             v               |
       |       [Call __get__]        |
       |       [Return result]       |
       |                             v
       +--------------------> [Check obj.__dict__]
                                     |
                             +-------+-------+
                             | Found         | Not Found
                             v               |
                       [Return dict val]     v
                                  [Is it a Non-Data Descriptor?]
                                  (has __get__ but no __set__)
                                             |
                                     +-------+-------+
                                     | Yes           | No
                                     v               v
                               [Call __get__]   [Is it class attribute?]
                               [Return result]       |
                                             +-------+-------+
                                             | Yes           | No
                                             v               v
                                       [Return value]   [Raise AttributeError]
                                                             |
                                                        [Call __getattr__]
\end{lstlisting}

\hypertarget{c-level-descriptor-specifications-pygetsetdef}{%
\subsubsection{\texorpdfstring{3. C-Level descriptor specifications:
\texttt{PyGetSetDef}}{3. C-Level descriptor specifications: PyGetSetDef}}\label{c-level-descriptor-specifications-pygetsetdef}}

In C extensions, properties are frequently defined using the
\passthrough{\lstinline!PyGetSetDef!} struct in the class's
\passthrough{\lstinline!tp\_getset!} slot:

\begin{lstlisting}[language=C]
typedef struct PyGetSetDef {
    const char *name;
    getter get;             /* C function: PyObject *(*getter)(PyObject *, void *) */
    setter set;             /* C function: int (*setter)(PyObject *, PyObject *, void *) */
    const char *doc;
    void *closure;          /* Context pointer passed to getter/setter */
} PyGetSetDef;
\end{lstlisting}

\hypertarget{practical-implementation-of-core-decorators}{%
\subsubsection{4. Practical Implementation of Core
Decorators}\label{practical-implementation-of-core-decorators}}

Here is how decorators leverage descriptors to modify binding behavior:
* \textbf{Bound Methods}: Python functions are non-data descriptors.
When accessed via \passthrough{\lstinline!obj.method!},
\passthrough{\lstinline!FunctionType.\_\_get\_\_(func, obj, Class)!} is
called. It returns a \passthrough{\lstinline!PyMethod!} object wrapping
the function and the instance:
\[\text{obj.method} \equiv \text{Class.method.\_\_get\_\_}(obj, \text{Class})\]
* \textbf{\passthrough{\lstinline!@classmethod!}}: Implements
\passthrough{\lstinline!\_\_get\_\_(self, instance, owner)!}. When
accessed, it ignores the \passthrough{\lstinline!instance!} argument and
returns a bound method wrapping the function and the
\passthrough{\lstinline!owner!} (the class object itself). *
\textbf{\passthrough{\lstinline!@staticmethod!}}: Implements
\passthrough{\lstinline!\_\_get\_\_(self, instance, owner)!}. It returns
the underlying raw function object directly, bypassing any method
binding. * \textbf{\passthrough{\lstinline!@property!}}: A data
descriptor that wraps getter, setter, and deleter functions: ```python
class CustomProperty(object): def \textbf{init}(self, fget, fset=None):
self.fget = fget self.fset = fset

\begin{lstlisting}
  def __get__(self, instance, owner):
      if instance is None:
          return self
      return self.fget(instance)

  def __set__(self, instance, value):
      if self.fset is None:
          raise AttributeError("can't set attribute")
      self.fset(instance, value)
\end{lstlisting}

\begin{lstlisting}
---

### 4.3 Method Resolution Order (MRO) & C3 Linearization

Method Resolution Order determines how Python traverses the inheritance tree during attribute search.

#### 1. The Classic DFLR Algorithm & The Diamond Problem
Classic classes (pre-Python 2.2) resolved attributes using a **Depth-First, Left-to-Right (DFLR)** tree traversal.
In a diamond inheritance hierarchy:
\end{lstlisting}

\begin{lstlisting}
 A
/ \
\end{lstlisting}

B C ~/ D

\begin{lstlisting}
The declaration `class D(B, C)` inherits from `B` and `C`, both of which inherit from `A`. 
* DFLR path: `[D, B, A, object, C, A, object]`.
* Removing duplicates keeping only the **first** occurrence gives: `[D, B, A, object, C]`.
* **The Failure**: If `A` defines a method `method()` and `C` overrides it, calling `D().method()` resolves to `A.method()` rather than `C.method()`. This violates the principle that specialized subclasses (`C`) should override generic ancestors (`A`).

#### 2. Python 2.2 MRO: The Last Occurrence Rule & Monotonicity Failure
To fix this, Python 2.2 introduced a new-style MRO calculation:
1. Perform DFLR traversal of the class and all its ancestors.
2. Remove all duplicates except the **last** occurrence.

For the diamond hierarchy `D(B, C)`:
* Traversal: `[D, B, A, object, C, A, object]`.
* Keeping only the **last** occurrence: `[D, B, C, A, object]`.
While this solved the diamond problem, the algorithm was **non-monotonic**.
* **Monotonicity**: If class $X$ precedes $Y$ in the MRO of class $P$, then $X$ must precede $Y$ in the MRO of any subclass $S$ derived from $P$.

##### Samuele Pedroni's Monotonicity Violation Example:
Consider this class configuration in Python 2.2:
```python
class A(object): pass
class B(object): pass
class C(object): pass
class D(object): pass
class E(object): pass

class K1(A, B, C): pass
class K2(D, B, E): pass
class K3(D, A): pass

class Z(K1, K2, K3): pass
\end{lstlisting}

Let's calculate the Python 2.2 MRO for \passthrough{\lstinline!K1!} and
\passthrough{\lstinline!K2!}: * \passthrough{\lstinline!K1!} raw:
\passthrough{\lstinline![K1, A, object, B, object, C, object]!}. Keeping
last: \passthrough{\lstinline![K1, A, B, C, object]!}. (Here, \(A\)
precedes \(B\)). * \passthrough{\lstinline!K2!} raw:
\passthrough{\lstinline![K2, D, object, B, object, E, object]!}. Keeping
last: \passthrough{\lstinline![K2, D, B, E, object]!}. (Here, \(B\)
precedes \(E\)). * \passthrough{\lstinline!K3!} raw:
\passthrough{\lstinline![K3, D, object, A, object]!}. Keeping last:
\passthrough{\lstinline![K3, D, A, object]!}.

Now let's resolve \passthrough{\lstinline!Z(K1, K2, K3)!} under Python
2.2: * Raw DFLR list:
\passthrough{\lstinline![Z, K1, A, object, B, object, C, object, K2, D, object, B, object, E, object, K3, D, object, A, object]!}.
* Keeping only the last occurrence:
\passthrough{\lstinline![Z, K1, C, K2, K3, D, B, E, A, object]!}. *
\textbf{The Monotonicity Failure}: * In the parent class
\passthrough{\lstinline!K1!}'s MRO, \(A\) preceded \(B\). * In the child
class \passthrough{\lstinline!Z!}'s MRO, \(B\) precedes \(A\)
(\passthrough{\lstinline!... D, B, E, A ...!}). * The child class has
reversed the relative order of \(A\) and \(B\) established in the
parent, violating monotonicity.

\hypertarget{python-2.3-mro-the-c3-linearization-algorithm}{%
\subsubsection{3. Python 2.3+ MRO: The C3 Linearization
Algorithm}\label{python-2.3-mro-the-c3-linearization-algorithm}}

To guarantee monotonicity and local precedence ordering, Python 2.3
adopted \textbf{C3 Linearization}.

\hypertarget{mathematical-formulation}{%
\paragraph{Mathematical Formulation:}\label{mathematical-formulation}}

Let \(L(C)\) be the linearization (MRO) of class \(C\). For a class
\(C\) inheriting from direct parents \(B_1, B_2, \dots, B_N\):
\[L(C) = [C] + \text{merge}\left(L(B_1), L(B_2), \dots, L(B_N), [B_1, B_2, \dots, B_N]\right)\]

Where \(L(\text{object}) = [\text{object}]\).

\hypertarget{the-merge-operation}{%
\paragraph{The Merge Operation:}\label{the-merge-operation}}

\begin{enumerate}
\def\labelenumi{\arabic{enumi}.}
\tightlist
\item
  Examine the head (index 0) of the first list inside the merge block:
  \(H = L(B_1)[0]\).
\item
  If \(H\) does not appear in the \textbf{tail} (index 1 to the end) of
  any other list in the merge block, it is a \textbf{good head}.

  \begin{itemize}
  \tightlist
  \item
    Append \(H\) to the linearization of \(C\).
  \item
    Remove \(H\) from all lists in the merge block.
  \item
    Repeat the merge step.
  \end{itemize}
\item
  If \(H\) appears in the tail of any other list, it is not a good head.
  Move to the next list in the merge block and check its head.
\item
  If no candidate head can be selected across all lists, the merge is
  impossible. Python raises a \passthrough{\lstinline!TypeError!}.
\end{enumerate}

\hypertarget{step-by-step-mathematical-calculation-of-zk1-k2-k3}{%
\paragraph{\texorpdfstring{Step-by-Step Mathematical Calculation of
\texttt{Z(K1,\ K2,\ K3)}:}{Step-by-Step Mathematical Calculation of Z(K1, K2, K3):}}\label{step-by-step-mathematical-calculation-of-zk1-k2-k3}}

Let's resolve the MRO of class \passthrough{\lstinline!Z!} from the
Pedroni example using C3 Linearization.

We have the parent linearizations:
\[L(K1) = [K1, A, B, C, \text{object}]\]
\[L(K2) = [K2, D, B, E, \text{object}]\]
\[L(K3) = [K3, D, A, \text{object}]\]

Now calculate \(L(Z)\):
\[L(Z) = [Z] + \text{merge}\left(L(K1), L(K2), L(K3), [K1, K2, K3]\right)\]
\[L(Z) = [Z] + \text{merge}\left([K1, A, B, C, \text{obj}], [K2, D, B, E, \text{obj}], [K3, D, A, \text{obj}], [K1, K2, K3]\right)\]

\textbf{Step 1}: Check head of first list: \passthrough{\lstinline!K1!}.
* Does \passthrough{\lstinline!K1!} appear in the tail of
\([K2, D, B, E, \text{obj}]\), \([K3, D, A, \text{obj}]\), or
\([K1, K2, K3]\)? No. * Extract \passthrough{\lstinline!K1!}:
\[L(Z) = [Z, K1] + \text{merge}\left([A, B, C, \text{obj}], [K2, D, B, E, \text{obj}], [K3, D, A, \text{obj}], [K2, K3]\right)\]

\textbf{Step 2}: Check head of first list: \passthrough{\lstinline!A!}.
* Does \passthrough{\lstinline!A!} appear in the tail of other lists?
Yes, it appears in the tail of \([K3, D, A, \text{obj}]\). Skip
\passthrough{\lstinline!A!}. * Move to the next list head:
\passthrough{\lstinline!K2!}. * Does \passthrough{\lstinline!K2!} appear
in the tail of other lists? No (it only appears at the head of the last
list \([K2, K3]\)). * Extract \passthrough{\lstinline!K2!}:
\[L(Z) = [Z, K1, K2] + \text{merge}\left([A, B, C, \text{obj}], [D, B, E, \text{obj}], [K3, D, A, \text{obj}], [K3]\right)\]

\textbf{Step 3}: Check head of first list: \passthrough{\lstinline!A!}.
* Does \passthrough{\lstinline!A!} appear in the tail of other lists?
Yes, in the tail of \([K3, D, A, \text{obj}]\). Skip. * Check head of
second list: \passthrough{\lstinline!D!}. * Does
\passthrough{\lstinline!D!} appear in the tail of other lists? Yes, in
the tail of \([K3, D, A, \text{obj}]\). Skip. * Check head of third
list: \passthrough{\lstinline!K3!}. * Does \passthrough{\lstinline!K3!}
appear in the tail of other lists? No. * Extract
\passthrough{\lstinline!K3!}:
\[L(Z) = [Z, K1, K2, K3] + \text{merge}\left([A, B, C, \text{obj}], [D, B, E, \text{obj}], [D, A, \text{obj}]\right)\]

\textbf{Step 4}: Check head of first list: \passthrough{\lstinline!A!}.
* Does \passthrough{\lstinline!A!} appear in the tail of other lists?
Yes, in the tail of \([D, A, \text{obj}]\). Skip. * Check head of second
list: \passthrough{\lstinline!D!}. * Does \passthrough{\lstinline!D!}
appear in the tail of other lists? No (it is at the head of
\([D, A, \text{obj}]\)). * Extract \passthrough{\lstinline!D!}:
\[L(Z) = [Z, K1, K2, K3, D] + \text{merge}\left([A, B, C, \text{obj}], [B, E, \text{obj}], [A, \text{obj}]\right)\]

\textbf{Step 5}: Check head of first list: \passthrough{\lstinline!A!}.
* Does \passthrough{\lstinline!A!} appear in the tail of other lists? No
(it is at the head of \([A, \text{obj}]\)). * Extract
\passthrough{\lstinline!A!}:
\[L(Z) = [Z, K1, K2, K3, D, A] + \text{merge}\left([B, C, \text{obj}], [B, E, \text{obj}], [\text{obj}]\right)\]

\textbf{Step 6}: Check head of first list: \passthrough{\lstinline!B!}.
* Does \passthrough{\lstinline!B!} appear in the tail of other lists?
No. * Extract \passthrough{\lstinline!B!}:
\[L(Z) = [Z, K1, K2, K3, D, A, B] + \text{merge}\left([C, \text{obj}], [E, \text{obj}], [\text{obj}]\right)\]

\textbf{Step 7}: Check head of first list: \passthrough{\lstinline!C!}.
* Does \passthrough{\lstinline!C!} appear in the tail of other lists?
No. * Extract \passthrough{\lstinline!C!}:
\[L(Z) = [Z, K1, K2, K3, D, A, B, C] + \text{merge}\left([\text{obj}], [E, \text{obj}], [\text{obj}]\right)\]

\textbf{Step 8}: Check head of first list:
\passthrough{\lstinline!obj!}. * Does \passthrough{\lstinline!obj!}
appear in the tail of other lists? Yes, in the tail of
\([E, \text{obj}]\). Skip. * Check head of second list:
\passthrough{\lstinline!E!}. * Does \passthrough{\lstinline!E!} appear
in the tail of other lists? No. * Extract \passthrough{\lstinline!E!}:
\[L(Z) = [Z, K1, K2, K3, D, A, B, C, E] + \text{merge}\left([\text{obj}], [\text{obj}], [\text{obj}]\right)\]

\textbf{Step 9}: Extract \passthrough{\lstinline!obj!}:
\[L(Z) = [Z, K1, K2, K3, D, A, B, C, E, \text{object}]\]

This calculation resolves the Method Resolution Order cleanly,
preserving both local precedence and global monotonicity across all
classes.

\hypertarget{instance-lifecycle-allocation-vs.-initialization}{%
\subsection{4.4 Instance Lifecycle: Allocation
vs.~Initialization}\label{instance-lifecycle-allocation-vs.-initialization}}

CPython splits object creation into two phases, controlled by different
type slots: 1. \textbf{Allocation
(\passthrough{\lstinline!\_\_new\_\_!})}: * Maps to the
\passthrough{\lstinline!tp\_new!} slot in
\passthrough{\lstinline!PyTypeObject!}. * This is a static method
responsible for allocating memory on the heap (using
\passthrough{\lstinline!PyType\_GenericNew()!}) and initializing the
object's \passthrough{\lstinline!PyObject!} headers
(\passthrough{\lstinline!ob\_refcnt!} and
\passthrough{\lstinline!ob\_type!}). * It must return a new instance of
the class (or subclass). 2. \textbf{Initialization
(\passthrough{\lstinline!\_\_init\_\_!})}: * Maps to the
\passthrough{\lstinline!tp\_init!} slot in
\passthrough{\lstinline!PyTypeObject!}. * This is an instance method
called immediately after \passthrough{\lstinline!\_\_new\_\_!} returns a
valid instance. It populates the instance dictionary
\passthrough{\lstinline!\_\_dict\_\_!} with fields.

\begin{center}\rule{0.5\linewidth}{0.5pt}\end{center}
